\documentclass[10pt,a4paper]{amsart}
\usepackage[margin=19mm]{geometry}
\usepackage{amsmath,amssymb,mathtools}
\usepackage[scr=rsfs,cal=euler]{mathalfa}
\usepackage{xfrac}
\usepackage[unicode,hidelinks,bookmarks=false]{hyperref}
\newcommand{\ie}{i.e.\ }
\newcommand{\cf}{cf.\ }
\newcommand{\resp}{respectively\ }
\newcommand{\Subsection}{\subsection}
\providecommand{\nobreakdash}{\nobreak\hbox{-}}
\setlength{\emergencystretch}{2em}
\multlinegap0pt
\usepackage[all]{xy}
\usepackage{tikz}
\usepackage{amssymb}
\usepackage{multirow}
\usepackage{mathtools}
\usepackage{xcolor}
\usepackage{booktabs}
\usepackage[shortlabels]{enumitem}
\usepackage{xfrac}
\newtheorem{lemma}{Lemma}
\title{An Arithmetic Topology viewpoint on Descent theory and Equivariant Categories}
\author{Miltiadis Karakikes, Sotiris Karanikolopoulos, Aristides Kontogeorgis, Dimitrios Noulas}
\date{Source: arXiv:2512.20551v1 (CC BY 4.0)}
\begin{document}
\maketitle

\noindent\emph{Source and licence.} An Arithmetic Topology viewpoint on Descent theory and Equivariant Categories. Version arXiv:2512.20551v1 (CC BY 4.0), distributed by arXiv under the Creative Commons Attribution 4.0 International licence, http://creativecommons.org/licenses/by/4.0/, as recorded in the arXiv licence metadata for this exact version and shown on the abstract page https://arxiv.org/abs/2512.20551v1. Full licence notice: \texttt{licenses/arxiv-2512.20551-CC-BY-4.0.txt}.\par\smallskip

\noindent\textit{Editorial scope.} Counted unit: the article's lemma lemma:conjNorm (source0.tex lines 1397--1408). The article's complete original proof of it is delivered with it (source0.tex lines 1409--1459, one \texttt{proof} environment of 51 lines). No mathematics is written, repaired or re-derived: the statement and the proof are the article's own lines, in the article's own notation, and no retained line is substituted or re-flowed.  No further source-local context is needed: every cross-reference of every retained line resolves inside the two delivered spans.  Nothing was omitted from any retained span, so the package carries no omission record.  No retained line contains a citation, so the wrapper carries no bibliography and no bibliography entry.  No retained line contains a figure environment, an image-inclusion command or drawing code, so no image is delivered and the package declares no asset dependency.  Editorial formatting only, in addition: an AMS \texttt{amsart} preamble that restates the article's own declarations and macros used by the retained lines, namely the environments \texttt{lemma} on separate counters, together with the standard packages those lines need.  The retained environments print in this document as Lemma 1 in order of appearance, whereas the article prints the same items with the numbering of its own full document.  The complete native source of the article as distributed in the arXiv e-print for this exact version is bundled unchanged as \texttt{sources/191579/source0.tex}.
\begin{lemma}
  \label{lemma:conjNorm}
Let $G$ be a group and $\bar{H}$ be a subgroup of $G$ of index $d$. Consider the natural left action of $G$ on cosets $\bar{H}, g_2 \bar{H}, \ldots, g_d \bar H$, which induces a representation
\[
  \Psi_{\bar H} \colon G \longrightarrow S_d. 
\]
For a conjugate subgroup $n\bar{H}n^{-1}$ of $\bar{H}$, the corresponding representation
\[
  \Psi_{n \bar H n^{-1}} \colon G \longrightarrow S_d
\]
is a conjugate of $\Psi_{\bar{H}}$ by an element $\varphi \in \mathrm{N}_{S_d}G$.
\end{lemma}

\begin{proof} By definition, for $g\in G$ we have 
\[
  g g_i =g_{\sigma(g)(i)} h, \ \text{for some } h \in \bar{H}
\]
and $\sigma(g)\in S_d$ is the permutation corresponding to $g$ under the injection $G\xhookrightarrow{}S_d$. 

For $n\bar{H}n^{-1}$, we can choose as coset representatives the elements
$n g_i n^{-1}$, for $i=1,\ldots,d$
 and the corresponding action of $ngn^{-1}$ is given by
\[
  n g  n^{-1} ng_in^{-1}  = n g_{\sigma(g)(i)} n^{-1} nhn^{-1}.
\]
This means that the element $ngn^{-1}$ acts on the cosets of $n\bar{H}n^{-1}$ in the same way that $ g $ acts on the cosets of $\bar{H}$ in $G$. This proves that $\Psi_{\bar H} (g) = \Psi_{n\bar H n^{-1}} (ngn^{-1})$, (up to conjugation by an element in $\mathrm{N}_{S_d}(G)$ for the specific choice of labeling of the cosets of $n\bar{H}n^{-1})$. Replace $g$ by $n^{-1}gn =\colon g^n $, so that $\Psi_{\bar H} (g^n) = \Psi_{n\bar H n^{-1}} (g)$.
The conjugation action by $n$ introduces an element $\tau\in S_d$, by writing
\[
  g_i^n = n^{-1}g_i n = g_{\tau(i)} h, \text{ for some element } h \in \bar{H}. 
\]
In the following diagram, we keep track of the coset representatives of $\bar{H}$ in the inner commutative diagram, whilst keeping track of the permutations in $S_d$ that arise after each operation on the representatives.
{\tiny
\[
\xymatrix{
  (1,\ldots,d) \ar[rrr]^{\tau}
    \ar[ddd]_{\sigma(g)} \ar@{~}[dr]  & & & (\tau(1),\ldots,\tau(d)) \ar[ddd]_{  \sigma(g^n) }
    \\
   &     \{g_1,\ldots,g_d\} 
  \ar[d]^{g_i \rightarrow g\cdot g_i}
  \ar[r]^{g_i\rightarrow g_i^n} & \{g_1^n,\ldots,g_d^n\} 
 \ar[d]^{g_i \rightarrow g^n\cdot g_i}
  \ar@{~}[ur] & 
  \\
    &     \{gg_1,\ldots,gg_d\} \ar[r]_{g_i\rightarrow g_i^n} & \{g^ng_1^n,\ldots,g^ng_d^n\} \ar@{~}[dr]
\\
(\sigma(g)(1),\ldots,\sigma(g)(d)) 
   \ar@{~}[ur]
   \ar[rrr]_\tau
& & & 
\substack{
  \{g_{\tau\sigma(g)(1)},\ldots,g_{\tau\sigma(g)(d)}\} \\
  = \\
   \{g_{\sigma(g^n)\tau(1)},\ldots,g_{\sigma(g^n)\tau(d)}\} 
}
}
\]
}

The commutativity of the above diagram, is interpreted by the relation
\[
  \sigma(g)\tau=\tau   \sigma(g^n),
\]
that is the element $n$ corresponds to  the element $\tau \in \mathrm{N}_{S_d}G$ as required.  
\end{proof}

\end{document}
